% GitHub Repo and Documentation: https://github.com/celiobjunior/resume-template
% Copyright © 2025 Celio B Junior. All rights reserved.
% 
% Licensed under the Apache License, Version 2.0 (the "License");
% you may not use this file except in compliance with the License.
% You may obtain a copy of the License at
%
%     http://www.apache.org/licenses/LICENSE-2.0
%
% This template follows best practices from README.md

\documentclass[a4paper,10pt]{article}

% --- PACOTES: BASE / IDIOMA ---
\usepackage[utf8]{inputenc}
\usepackage[T1]{fontenc}
\usepackage[portuguese]{babel}

% --- PACOTES: LAYOUT / TIPOGRAFIA ---
\usepackage{geometry}
\usepackage{parskip}
\usepackage{microtype}
\usepackage{enumitem}
\usepackage{titlesec}
\usepackage{array}
\usepackage{tabularx}

% --- PACOTES: LINKS / BOOKMARKS ---
\usepackage{hyperref}
\usepackage{bookmark}
\usepackage{xurl}

% --- CONFIGURAÇÃO: MARGENS / ESPAÇAMENTO ---
\geometry{top=1.5cm, bottom=1.5cm, left=1.5cm, right=1.5cm}
\setcounter{secnumdepth}{0}
\setlist[itemize]{
    leftmargin=0.75em,
    itemsep=0.2em,
    topsep=0.25em,
    parsep=0em,
    partopsep=0em
}

\pagestyle{empty}

% --- CONFIGURAÇÃO: METADADOS / LINKS ---
\hypersetup{
    pdftitle={CV Lucas Correia da Silva},
    pdfauthor={Lucas Correia da Silva},
    colorlinks=true,
    linkcolor=black,
    urlcolor=black,
    citecolor=black,
    bookmarksdepth=2
}

% --- CONFIGURAÇÃO: TÍTULOS ---
\titleformat{\section}
{\Large\bfseries}
{}
{0em}
{}
[\titlerule\vspace{0.5ex}]
\titleformat{\subsection}
{\normalsize\bfseries}
{}
{0em}
{}
\titlespacing*{\subsection}{0pt}{0.35em}{0.15em}

% Macro para entradas do currículo com bookmark no PDF
\newcounter{cventry}
\newcommand{\cventry}[4]{%
    \refstepcounter{cventry}%
    \phantomsection%
    \pdfbookmark[2]{#1}{cventry-\thecventry}%
    \noindent\begin{tabularx}{\textwidth}{@{}>{\raggedright\arraybackslash}X >{\raggedleft\arraybackslash}X@{}}
    \textbf{#1} & #2 \\
    \textit{#3} & \textit{#4} \\
    \end{tabularx}
}

% --- INÍCIO DO DOCUMENTO ---
\begin{document}

% --- CABEÇALHO ---
% FIX ATS: separadores em texto puro (| ao invés de \textbullet)
% evita concatenação de campos pelo parser
\begin{center}
    {\LARGE \textbf{Lucas Correia da Silva}}
    \\ [0.1cm]
    Vitória, ES
    { | }
    \href{mailto:lucasdevproject@gmail.com}{lucasdevproject@gmail.com}
    { | }
    \href{https://www.linkedin.com/in/olucascdev/}{linkedin.com/in/olucascdev}
    { | }
    \href{https://github.com/olucascdev}{github.com/olucascdev}
\end{center}

% --- SEÇÕES ---

\section{Objetivo}
Desenvolvedor Back-End/FullStack com experiência no desenvolvimento de sistemas web, automações, integrações e soluções orientadas a APIs. Atuação com PHP, Python e Go, utilizando frameworks modernos, Docker e arquitetura escalável para construção de aplicações robustas, organizadas e eficientes.

% ------

\section{Experiência}

\cventry{EVA Tecnologia}{Vitória, ES}{Desenvolvedor de Software}{Jan 2026 -- Atual}
\begin{itemize}
    \item Atuação no desenvolvimento de soluções de automação e integração de sistemas utilizando Python, Agno e ferramentas baseadas em IA.
    \item Criação de fluxos inteligentes com n8n para atendimento automatizado e otimização de processos internos.
    \item Utilização de Docker e arquitetura orientada a APIs para garantir escalabilidade, organização e eficiência dos serviços.
\end{itemize}

\cventry{JN Desenvolvimento de Sistemas LTDA}{Vitória, ES}{Desenvolvedor de Software Jr Full Stack}{Abr 2025 -- Out 2025}
\begin{itemize}
    \item Desenvolvimento de sistemas web sob medida para diferentes nichos empresariais utilizando Laravel e Filament.
    \item Implementação de arquitetura multitenancy para garantir escalabilidade e isolamento seguro de dados entre múltiplos clientes.
    \item Integração com APIs de terceiros, incluindo pagamentos, autenticação e serviços externos, promovendo interoperabilidade entre sistemas.
    \item Criação de dashboards administrativos, fluxos automatizados e interfaces dinâmicas com foco em performance e usabilidade.
    \item Estruturação de regras de negócio e otimização de processos internos, reduzindo retrabalho e aumentando a eficiência operacional.
\end{itemize}

\cventry{Perbras}{Vitória, ES}{Estagiário de TI}{Ago 2024 -- Abr 2025}
\begin{itemize}
    \item Desenvolvimento de soluções internas voltadas à automação de processos operacionais, contribuindo para o aumento da produtividade das equipes e redução de custos.
    \item Atuação em suporte técnico em sistemas industriais críticos, garantindo alta disponibilidade e rápida resolução de incidentes.
    \item Integração e otimização de fluxos de trabalho, melhorando a eficiência dos sistemas utilizados em campo.
    \item Criação de treinamentos imersivos em Realidade Virtual com C\# e Unity voltados à capacitação técnica e segurança operacional.
\end{itemize}

% ------

\section{Projetos}

% FIX ATS: datas adicionadas em todos os projetos para evitar confusão no parser de linha do tempo

\cventry{Illie}{Projeto Pessoal}{Assistente pessoal de conteúdo com IA}{2025 -- Atual}
\begin{itemize}
    \item Assistente pessoal de conteúdo que transforma estudos e projetos escritos no Notion em posts prontos para o LinkedIn, automaticamente, no estilo do usuário.
    \item Monitora páginas do Notion marcadas como prontas, reescreve o conteúdo como um post autêntico e armazena tudo em banco de dados para revisão via dashboard web.
    \item Tecnologias: Python, Agno, Kestra, PostgreSQL, Next.js e Docker.
\end{itemize}

\cventry{Edully -- Sistema de Gestão Educacional}{Projeto Pessoal / Sob Demanda}{Sistema de gestão educacional}{2024 -- Atual}
\begin{itemize}
    \item Sistema de gestão educacional desenvolvido para integrar processos acadêmicos, financeiros e administrativos.
    \item Possui múltiplos painéis por perfil de usuário e arquitetura escalável baseada em RBAC.
    \item Tecnologias: Laravel, Filament, PHP, Livewire, TailwindCSS, PostgreSQL, Docker e MinIO.
\end{itemize}

\cventry{Assistente de Qualificação de Leads com IA}{Desafio Técnico}{Sistema inteligente de qualificação de leads}{2025}
\begin{itemize}
    \item Sistema desenvolvido com agentes de IA para conduzir conversas naturais, coletar informações estratégicas e gerar relatórios estruturados com score de qualificação.
    \item Integra base de conhecimento, cálculo automático de orçamentos e armazenamento de histórico.
    \item Tecnologias: Python, Agno e SQLite.
\end{itemize}

\cventry{Fórum de Comunidades (Clone Reddit)}{Desafio Técnico}{Aplicação full stack inspirada no Reddit}{2024}
\begin{itemize}
    \item Aplicação com criação dinâmica de comunidades, publicação de conteúdos em Markdown, comentários e sistema de votos.
    \item Inclui painel administrativo para gerenciamento, com foco em modelagem de dados, arquitetura e organização de código.
    \item Tecnologias: PHP, Laravel, Filament, Blade, TailwindCSS, PostgreSQL e Docker.
\end{itemize}

\cventry{LunchTube}{Projeto Pessoal}{Extensão de navegador}{2024}
\begin{itemize}
    \item Extensão de navegador focada em otimizar momentos de pausa, eliminando o tempo gasto na escolha de vídeos no YouTube.
    \item Tecnologias: JavaScript, HTML e CSS.
\end{itemize}

\cventry{BugTracker}{Projeto Pessoal}{Sistema de monitoramento e gestão de erros}{2025}
\begin{itemize}
    \item Sistema de monitoramento e gestão de erros em atendimentos automatizados por IA.
    \item Permite análise de qualidade, priorização de falhas e acompanhamento de métricas como taxa de assertividade, volume de atendimentos e recorrência de erros.
    \item Possui visões separadas para equipe técnica e cliente.
    \item Tecnologias: Laravel, Filament, Docker e PostgreSQL.
\end{itemize}

% ------

\section{Habilidades}
\begin{itemize}
    \item \textbf{Back-end:} Go, Python, PHP
    \item \textbf{Frameworks:} Gin, Laravel, Filament, Flask, FastAPI, Agno
    \item \textbf{Banco de Dados:} PostgreSQL, MySQL, MongoDB, Redis
    \item \textbf{Infrastructure \& DevOps:} Docker, Linux Servers, VPS Deployment, CI/CD, GitHub Actions, Kestra
    \item \textbf{Front-end:} HTML, CSS, JavaScript, TailwindCSS, Next.js
    \item \textbf{Arquitetura:} RESTful APIs, Microservices, Multitenancy
    \item \textbf{Idiomas:} Português (Nativo), Inglês (B1)
\end{itemize}

% ------

\section{Educação}

\cventry{Centro Universitário Vale do Cricaré -- UNIVC}{Vitória, ES}{CTS em Análise e Desenvolvimento de Sistemas}{Atualmente -- 5º Período}
\begin{itemize}
    \item Participação ativa na Jornada Científica Semestral da Universidade, contribuindo para apresentações, discussões acadêmicas e projetos interdisciplinares.
\end{itemize}

\cventry{Universidade Federal do Espírito Santo -- UFES}{Espírito Santo, Brasil}{Ciência da Computação, Bacharelado (Incompleto)}{Ago 2023 -- Maio 2024}
\begin{itemize}
    \item Membro da empresa júnior Adapti na área de Design e Marketing.
    \item Vencedor na categoria de calouros do Campeonato de Programação (MAPCOM).
\end{itemize}

\end{document}
